\chapter{Phương pháp xây dựng Vietnamese NeoBERT}
\label{chap:method}

Chương này trình bày chi tiết phương pháp xây dựng mô hình Vietnamese NeoBERT
(ViNeoBERT). Trước hết, chúng tôi mô tả tổng quan quy trình ba giai đoạn cùng các mô hình
nguồn được sử dụng (Mục~\ref{sec:method-overview}). Tiếp theo, chúng tôi trình bày bước cắt
tỉa và đồng bộ từ vựng (Mục~\ref{sec:vocab}), quá trình trích xuất tập neo theo cặp dịch --
một cải tiến của đề tài so với SALT gốc (Mục~\ref{sec:anchor}), và phép chiếu SALT áp dụng
cho cả lớp nhúng đầu vào lẫn lớp đầu ra tách rời của NeoBERT (Mục~\ref{sec:salt-proj}).
Cuối cùng, chúng tôi mô tả giai đoạn căn chỉnh đóng băng (Mục~\ref{sec:freeze}) và giai
đoạn tiếp tục tiền huấn luyện thích nghi (Mục~\ref{sec:cpt}). Các giá trị siêu tham số cụ
thể được trình bày trong Chương~\ref{chap:experiment}.

% ======================================================================
\section{Tổng quan phương pháp}
\label{sec:method-overview}

\subsection{Đặt vấn đề}
\label{subsec:problem}

Như đã phân tích trong Chương~\ref{chap:related}, việc xây dựng một Encoder hiện đại cho
tiếng Việt từ NeoBERT đối mặt đồng thời ba thách thức: (1)~NeoBERT chỉ được tiền huấn luyện
trên ngữ liệu tiếng Anh, nên toàn bộ lớp nhúng của nó phản ánh phân phối từ vựng Anh ngữ;
(2)~tiếp tục tiền huấn luyện với một lớp nhúng \textit{khởi tạo ngẫu nhiên} dễ dẫn tới hiện
tượng quên lãng thảm họa; và (3)~huấn luyện lại từ đầu trên quy mô hàng nghìn tỷ token là
bất khả thi về tài nguyên. Phương pháp của đề tài giải quyết đồng thời ba thách thức này
bằng cách \textit{khởi tạo có thông tin} lớp nhúng tiếng Việt thông qua kỹ thuật SALT, sau
đó tinh chỉnh mô hình trên ngữ liệu tiếng Việt với chi phí khiêm tốn.

\subsection{Mô hình nguồn và mô hình hiến tặng}
\label{subsec:sources}

Phương pháp sử dụng hai mô hình tiền huấn luyện:
\begin{itemize}
  \item \textbf{NeoBERT}~\cite{lebreton2025neobert} đóng vai trò \textit{kiến trúc đích}:
        toàn bộ phần thân Transformer (28 lớp với RoPE, SwiGLU, Pre-RMSNorm) cùng không gian
        biểu diễn tiếng Anh được giữ nguyên và dùng làm hệ tọa độ đích cho phép chiếu. Một
        đặc điểm quan trọng là NeoBERT sử dụng \textit{đầu giải mã tách rời} (untied head):
        lớp đầu ra của tác vụ MLM là một ánh xạ tuyến tính riêng $\,y = W_{\mathrm{dec}}\,h
        + b_{\mathrm{dec}}\,$ với ma trận trọng số $W_{\mathrm{dec}}$ và véc-tơ độ chệch
        $b_{\mathrm{dec}}$ \textit{độc lập} với ma trận nhúng đầu vào $E$. Do đó, việc thích
        nghi sang tiếng Việt đòi hỏi khởi tạo \textit{cả hai} ma trận token: ma trận nhúng
        đầu vào $E$ và ma trận đầu ra $W_{\mathrm{dec}}$.
  \item \textbf{ViDeBERTa}~\cite{lai2023videberta} (bản \texttt{base}) đóng vai trò
        \textit{mô hình hiến tặng} (donor): cung cấp không gian nhúng tiếng Việt chất lượng
        cao làm nguồn để tái sử dụng. ViDeBERTa dùng bộ tách từ SentencePiece theo thuật
        toán Unigram với kích thước từ vựng lớn (khoảng 128 nghìn token), trong đó mỗi đơn
        vị từ vựng đi kèm một điểm log-xác suất.
\end{itemize}

\subsection{Quy trình ba giai đoạn}
\label{subsec:pipeline}

Mô hình ViNeoBERT được xây dựng qua ba giai đoạn nối tiếp, minh họa trong
Hình~\ref{fig:pipeline}:
\begin{enumerate}
  \item \textbf{Khởi tạo SALT}: xây dựng tập từ vựng tiếng Việt, trích xuất tập neo, rồi
        chiếu không gian nhúng của ViDeBERTa sang không gian NeoBERT để khởi tạo cả ma trận
        nhúng đầu vào lẫn ma trận đầu ra; độ chệch của đầu ra được khởi tạo theo tần suất từ
        đơn (Mục~\ref{sec:vocab}--\ref{sec:salt-proj}). Giai đoạn này \textit{không} cập
        nhật phần thân Transformer.
  \item \textbf{Căn chỉnh đóng băng}: đóng băng toàn bộ phần thân, chỉ huấn luyện ma trận
        nhúng và ma trận đầu ra trên một lượng nhỏ ngữ liệu tiếng Việt, nhằm để hai ma trận
        vừa khởi tạo căn chỉnh với nhau và với phần thân (Mục~\ref{sec:freeze}).
  \item \textbf{Tiếp tục tiền huấn luyện thích nghi}: tinh chỉnh toàn bộ mô hình trên ngữ
        liệu tiếng Việt quy mô lớn với mục tiêu MLM theo lịch tốc độ học WSD
        (Mục~\ref{sec:cpt}).
\end{enumerate}

\begin{figure}[ht]
  \centering
  \begin{tikzpicture}[
    node distance=7mm and 12mm,
    stage/.style={draw, rounded corners, minimum width=38mm, minimum height=13mm, align=center, font=\small, fill=blue!7},
    src/.style={draw, rounded corners, minimum width=30mm, minimum height=9mm, align=center, font=\footnotesize, fill=gray!10},
    >={Latex[length=2mm]}
  ]
    \node[src] (neo)  {NeoBERT\\(kiến trúc + không gian đích)};
    \node[src, below=of neo] (vide) {ViDeBERTa\\(mô hình hiến tặng)};
    \node[stage, right=of $(neo)!0.5!(vide)$] (s1) {Giai đoạn 1\\Khởi tạo SALT};
    \node[stage, right=of s1] (s2) {Giai đoạn 2\\Căn chỉnh đóng băng};
    \node[stage, right=of s2] (s3) {Giai đoạn 3\\CPT (WSD)};
    \node[src, right=of s3] (out) {ViNeoBERT};
    \draw[->] (neo.east)  -- (s1.west);
    \draw[->] (vide.east) -- (s1.west);
    \draw[->] (s1) -- (s2); \draw[->] (s2) -- (s3); \draw[->] (s3) -- (out);
    \node[font=\scriptsize, below=1mm of s1] {thân đóng băng};
    \node[font=\scriptsize, below=1mm of s2] {thân đóng băng};
    \node[font=\scriptsize, below=1mm of s3] {toàn bộ tham số};
  \end{tikzpicture}
  \caption{Quy trình ba giai đoạn xây dựng ViNeoBERT: khởi tạo lớp nhúng và lớp đầu ra bằng
    SALT từ ViDeBERTa sang không gian NeoBERT, căn chỉnh đóng băng phần thân, rồi tiếp tục
    tiền huấn luyện thích nghi trên ngữ liệu tiếng Việt.}
  \label{fig:pipeline}
\end{figure}

% ======================================================================
\section{Cắt tỉa và đồng bộ từ vựng}
\label{sec:vocab}

Phép chiếu SALT yêu cầu lớp nhúng nguồn và đích có cùng kích thước từ vựng. Vì từ vựng của
ViDeBERTa lớn hơn nhiều so với 30.522 token của NeoBERT, cần một bước \textit{cắt tỉa từ
vựng} (vocabulary pruning) để đồng bộ kích thước. Quá trình này giữ lại toàn bộ các token
đặc biệt (như \texttt{[CLS]}, \texttt{[SEP]}, \texttt{[MASK]}, \texttt{[PAD]}), sau đó chọn
bổ sung các token thông thường có \textit{điểm log-xác suất Unigram cao nhất} -- tức các đơn
vị từ vựng phổ biến và đáng tin cậy nhất -- cho tới khi đạt đúng 30.522 token.

Gọi $\mathcal{V}_s$ là tập token đặc biệt và $\mathcal{V}_n$ là tập token thông thường của
ViDeBERTa, với điểm Unigram $\mathrm{score}(\cdot)$. Tập từ vựng tiếng Việt sau cắt tỉa
$\mathcal{V}_t$ được xác định như Công thức~\eqref{eq:prune}:
\begin{equation}
  \label{eq:prune}
  \mathcal{V}_t = \mathcal{V}_s \;\cup\; \mathrm{TopK}\big(\mathcal{V}_n,\, \mathrm{score},\,
  K\big), \qquad K = 30{,}522 - |\mathcal{V}_s|
\end{equation}
Trong đó $\mathrm{TopK}(\mathcal{V}_n, \mathrm{score}, K)$ trả về $K$ token có điểm Unigram
cao nhất. Bộ tách từ được ghi lại với ánh xạ chỉ số mới, đồng thời lưu một ánh xạ
$\pi:\, \text{id mới} \rightarrow \text{id gốc trong ViDeBERTa}$ để truy xuất đúng hàng nhúng
hiến tặng ở các bước sau. Phép so khớp chuỗi piece là chính xác vì các đơn vị Unigram là duy
nhất trong từ vựng.

% ======================================================================
\section{Trích xuất tập neo theo cặp dịch}
\label{sec:anchor}

\subsection{Vai trò của tập neo}
\label{subsec:anchor-role}

\textit{Tập neo} (anchor set) $\mathcal{A}$ là tập các cặp token $(t_{\mathrm{vi}},
t_{\mathrm{neo}})$ có tương ứng ngữ nghĩa rõ ràng giữa từ vựng tiếng Việt và từ vựng NeoBERT.
Mỗi cặp neo cung cấp một điểm \textit{đã biết} của ánh xạ giữa hai không gian nhúng: nhúng
NeoBERT của $t_{\mathrm{neo}}$ chính là đích mà nhúng ViDeBERTa của $t_{\mathrm{vi}}$ cần
được ánh xạ tới. Do đó, chất lượng và độ phủ của tập neo quyết định trực tiếp chất lượng của
phép chiếu SALT.

Trong nghiên cứu gốc, SALT~\cite{lee2025salt} cũng như FOCUS~\cite{dobler2023focus} chủ yếu
khai thác các token \textit{trùng lặp bề mặt} (cùng dạng chuỗi) giữa hai từ vựng làm neo. Đối
với cặp ngôn ngữ Anh--Việt vốn khác biệt lớn về hình thái, số token trùng lặp bề mặt là khá
hạn chế. Vì vậy, đề tài đề xuất \textit{mở rộng tập neo bằng các cặp dịch} Anh--Việt được
khai thác tự động, nhằm làm giàu và phủ rộng tập neo. Đây là một trong những đóng góp chính
của đề tài.

\subsection{Lọc từ đầy đủ và hợp lệ}
\label{subsec:anchor-filter}

Để bảo đảm các cặp neo mang ngữ nghĩa trong sạch, trước khi khai thác cặp dịch, chúng tôi áp
dụng hai bộ lọc. Thứ nhất, chỉ giữ lại các \textit{token là từ đầy đủ} (word-initial), loại
bỏ các mảnh từ phụ (subword) vốn không mang nghĩa độc lập. Thứ hai, đối với token tiếng Việt,
chúng tôi áp dụng một bộ \textit{lọc chính tả tiếng Việt} nghiêm ngặt: loại các token chứa ký
tự không bản địa (\textit{f, j, w, z}), giới hạn độ dài, yêu cầu đúng một cụm nguyên âm liền
mạch, và tuân thủ các quy tắc chính tả (chẳng hạn \textit{k, gh, ngh} chỉ đứng trước
\textit{i, e, ê, y}). Hai bộ lọc này bảo đảm chỉ những từ tiếng Việt hợp lệ và đơn nghĩa mới
được xét làm neo.

\subsection{Khai thác cặp dịch ba tầng}
\label{subsec:anchor-mining}

Trên tập token đã lọc, chúng tôi khai thác cặp dịch theo ba tầng bổ sung cho nhau:
\begin{itemize}
  \item \textbf{Tầng trùng lặp bề mặt và số}: các từ có cùng dạng chuỗi ở cả hai từ vựng
        (được xác minh bằng dịch máy MarianMT~\cite{junczys2018marian}) và các token số chung.
  \item \textbf{Tầng dịch ngược một từ}: với mỗi từ tiếng Việt, chúng tôi dịch sang tiếng
        Anh rồi \textit{dịch ngược} trở lại tiếng Việt; cặp neo chỉ được chấp nhận nếu kết
        quả dịch ngược trùng khớp với từ gốc và bản dịch tiếng Anh là một từ đơn thuộc từ vựng
        NeoBERT.
  \item \textbf{Tầng từ ghép}: các từ ghép tiếng Việt (đa âm tiết) được dịch trực tiếp, chỉ
        giữ lại khi bản dịch là một từ đơn thuộc từ vựng NeoBERT.
\end{itemize}

Cơ sở khoa học của phép kiểm tra \textit{dịch ngược} là: một từ đa nghĩa thường không dịch
ngược về chính nó một cách ổn định, do bản dịch trung gian có thể tương ứng với một nghĩa
khác. Vì vậy, ràng buộc dịch ngược trùng khớp tự nhiên loại bỏ phần lớn các từ đa nghĩa, chỉ
giữ lại các cặp tương ứng một--một đáng tin cậy. Ràng buộc ``bản dịch là một từ đơn'' bảo đảm
mỗi cặp neo tương ứng đúng một token NeoBERT. Toàn bộ tập neo thu được hợp nhất từ ba tầng và
được dùng thống nhất cho cả lớp nhúng đầu vào lẫn lớp đầu ra ở Mục~\ref{sec:salt-proj}.

% ======================================================================
\section{Phép chiếu SALT cho lớp nhúng và lớp đầu ra}
\label{sec:salt-proj}

Mục này trình bày phép chiếu SALT -- bước khởi tạo cốt lõi. Với mỗi token tiếng Việt
\textit{không thuộc} tập neo, chúng tôi xây dựng nhúng của nó bằng một ánh xạ tuyến tính cục
bộ học từ các neo lân cận về mặt ngữ nghĩa. Các token neo và token đặc biệt được sao chép
trực tiếp hàng tương ứng từ NeoBERT.

\subsection{Lựa chọn neo bằng Sparsemax}
\label{subsec:anchor-select}

Để xác định các neo liên quan tới một token đích $t$, chúng tôi đo độ tương đồng ngữ nghĩa
trong không gian nhúng từ tĩnh fastText~\cite{bojanowski2017fasttext, grave2018fasttext}. Gọi
$\mathbf{f}_t$ là véc-tơ fastText của $t$ và $\{\mathbf{f}_a\}_{a\in\mathcal{A}}$ là véc-tơ
fastText của các neo. Độ tương đồng cosine và trọng số neo được tính như
Công thức~\eqref{eq:sparsemax-sel}:
\begin{equation}
  \label{eq:sparsemax-sel}
  s_{t,a} = \frac{\mathbf{f}_t^\top \mathbf{f}_a}{\lVert \mathbf{f}_t\rVert\, \lVert \mathbf{f}_a\rVert},
  \qquad
  \mathbf{w}_t = \mathrm{sparsemax}\big(\mathbf{s}_t\big),
  \qquad
  \mathcal{S}_t = \{\, a : w_{t,a} > 0 \,\}
\end{equation}
Trong đó $\mathrm{sparsemax}(\cdot)$~\cite{martins2016sparsemax} là phép chiếu lên đơn hình
xác suất, cho kết quả \textit{thưa} -- chỉ một tập con nhỏ $\mathcal{S}_t$ các neo liên quan
nhất nhận trọng số khác 0 (xem Mục~\ref{subsec:init-methods}). Việc dùng Sparsemax thay cho
softmax đặc biệt phù hợp với cấu trúc \textit{bất đẳng hướng} của không gian biểu diễn
(Mục~\ref{sec:geometry}): vì trong không gian bất đẳng hướng mọi cặp từ đều có độ tương đồng
cosine khá cao, một phân phối dày đặc sẽ pha trộn nhiễu, trong khi hỗ trợ thưa giữ lại đúng
các neo thực sự gần nghĩa.

\subsection{Ánh xạ tuyến tính cục bộ theo từng token}
\label{subsec:local-map}

Như đã lập luận ở Mục~\ref{sec:geometry}, do hai không gian nhúng là bất đẳng hướng và không
đẳng cấu, một ánh xạ tuyến tính toàn cục duy nhất là không đủ. Thay vào đó, SALT học một
\textit{ánh xạ tuyến tính riêng cho từng token} từ tập neo được chọn. Gọi $A_t^{\mathrm{src}}
\in \mathbb{R}^{k\times h}$ là ma trận xếp chồng các nhúng \textit{hiến tặng} (ViDeBERTa) của
$k = |\mathcal{S}_t|$ neo trong $\mathcal{S}_t$, và $A_t^{\mathrm{tgt}} \in \mathbb{R}^{k\times
h}$ là ma trận các hàng \textit{đích} (NeoBERT) tương ứng. Khi số neo đủ lớn
($k \ge k_0$), ánh xạ cục bộ $X_t$ và nhúng kết quả $\mathbf{e}_t$ được xác định như
Công thức~\eqref{eq:local-map}:
\begin{equation}
  \label{eq:local-map}
  X_t = \big(A_t^{\mathrm{src}}\big)^{+} A_t^{\mathrm{tgt}}, \qquad
  \mathbf{e}_t = \mathbf{d}_t\, X_t
\end{equation}
Trong đó $(\cdot)^{+}$ là \textit{giả nghịch đảo} Moore--Penrose (nghiệm bình phương tối
thiểu chuẩn nhỏ nhất), $\mathbf{d}_t \in \mathbb{R}^{h}$ là nhúng hiến tặng của token $t$, và
$h$ là số chiều ẩn. Ánh xạ $X_t$ học quan hệ tuyến tính giữa hai không gian \textit{cục bộ
quanh token $t$}, dựa trên các neo gần nghĩa nhất với nó.

Khi số neo không đủ ($k < k_0$), giả nghịch đảo trở nên kém ổn định; khi đó chúng tôi thoái
lui về \textit{trung bình có trọng số Sparsemax} của các hàng đích, như
Công thức~\eqref{eq:wavg}:
\begin{equation}
  \label{eq:wavg}
  \mathbf{e}_t = \frac{\sum_{a\in\mathcal{S}_t} w_{t,a}\, \mathbf{a}^{\mathrm{tgt}}}
                      {\sum_{a\in\mathcal{S}_t} w_{t,a}}
\end{equation}
Trường hợp token không có véc-tơ fastText hợp lệ, nhúng được khởi tạo ngẫu nhiên theo phân
phối chuẩn khớp với thống kê của NeoBERT.

\subsection{Áp dụng cho cả lớp nhúng đầu vào và lớp đầu ra}
\label{subsec:both-matrices}

Đây là điểm khác biệt then chốt so với SALT gốc. Vì NeoBERT dùng đầu giải mã \textit{tách
rời}, ma trận đầu ra $W_{\mathrm{dec}}$ là một ma trận độc lập với hình học riêng -- và theo
kết quả về hiện tượng thoái hóa biểu diễn (Mục~\ref{sec:geometry}), nó \textit{không} thể
được khởi tạo bằng cách sao chép ma trận nhúng. Do đó, chúng tôi áp dụng \textit{cùng một quy
trình chiếu SALT theo từng token} cho cả hai ma trận, chỉ khác ở tập hàng đích:
\begin{itemize}
  \item Với \textbf{ma trận nhúng đầu vào} $E$: hàng đích $A_t^{\mathrm{tgt}}$ lấy từ các
        \textit{hàng nhúng} NeoBERT của các neo.
  \item Với \textbf{ma trận đầu ra} $W_{\mathrm{dec}}$: hàng đích $A_t^{\mathrm{tgt}}$ lấy từ
        các \textit{hàng đầu giải mã} NeoBERT của các neo.
\end{itemize}
Nhờ vậy, mỗi ma trận được khởi tạo trong đúng không gian của nó, phản ánh trung thực quan hệ
ngữ nghĩa mà các neo thiết lập trong từng không gian.

\subsection{Hiệu chỉnh chuẩn của lớp nhúng}
\label{subsec:norm-calib}

Như đã lưu ý ở Mục~\ref{subsec:rmsnorm}, NeoBERT dùng Pre-RMSNorm vốn chuẩn hóa theo độ lớn
của véc-tơ; do đó độ lớn (norm) của mỗi hàng nhúng ảnh hưởng trực tiếp đến hành vi mô hình.
Các nhúng sinh ra từ phép chiếu SALT thường có độ lớn lệch so với phân phối của NeoBERT. Để
tránh gây bất ổn định, chúng tôi hiệu chỉnh độ lớn của mỗi hàng \textit{không phải neo} về
độ lớn trung bình của NeoBERT trong khi giữ nguyên hướng, như Công thức~\eqref{eq:norm-calib}:
\begin{equation}
  \label{eq:norm-calib}
  \mathbf{e}_t \;\leftarrow\; \mathbf{e}_t \cdot \frac{\mu_{\mathrm{neo}}}{\lVert \mathbf{e}_t\rVert},
  \qquad
  \mu_{\mathrm{neo}} = \frac{1}{|\mathcal{V}_{\mathrm{neo}}|}\sum_{v} \lVert E^{\mathrm{neo}}_v\rVert
\end{equation}
Trong đó $\mu_{\mathrm{neo}}$ là độ lớn trung bình của các hàng nhúng NeoBERT. Các hàng neo
và token đặc biệt được giữ nguyên để bảo toàn độ trung thực với NeoBERT.

\subsection{Khởi tạo độ chệch đầu ra theo tần suất}
\label{subsec:freq-bias}

Lớp đầu ra của NeoBERT có một véc-tơ độ chệch $b_{\mathrm{dec}}$ mã hóa \textit{tiên nghiệm
tần suất} của token. Các phương pháp khởi tạo trước đây thường để độ chệch này bằng 0, khiến
mô hình ở bước đầu phải lãng phí công sức học lại phân phối tần suất biên. Đề tài khởi tạo
$b_{\mathrm{dec}}$ theo \textit{log-tần-suất từ đơn} (unigram log-frequency) của tiếng Việt,
như Công thức~\eqref{eq:freq-bias}:
\begin{equation}
  \label{eq:freq-bias}
  b_{\mathrm{dec},v} = \log p_v, \qquad
  p_v = \frac{c_v + 1}{\sum_{u\in\mathcal{V}_t} c_u + |\mathcal{V}_t|}
\end{equation}
Trong đó $c_v$ là số lần xuất hiện của token $v$ trong ngữ liệu tiếng Việt (đếm qua bộ tách
từ đã cắt tỉa), và phép cộng 1 ở tử số là làm trơn Laplace. Nhờ cách khởi tạo này, ngay ở
bước huấn luyện đầu tiên mô hình đã ``đoán'' đúng phân phối tần suất từ đơn thay vì phân phối
đều, giúp hàm mất mát khởi điểm xấp xỉ \textit{entropy từ đơn} của tiếng Việt -- thấp hơn
đáng kể so với mức ngẫu nhiên $\log|\mathcal{V}_t|$. Đây là một cải tiến đơn giản nhưng hiệu
quả, phối hợp ăn khớp với việc khởi tạo trọng số đầu ra bằng SALT.

% ======================================================================
\section{Giai đoạn căn chỉnh đóng băng}
\label{sec:freeze}

Sau bước khởi tạo, ma trận nhúng và ma trận đầu ra của tiếng Việt là \textit{mới} đối với
phần thân Transformer đã được huấn luyện trên tiếng Anh. Nếu lập tức tinh chỉnh toàn bộ mô
hình, phần thân có thể bị kéo theo để thích ứng với hai ma trận mới còn chưa ổn định, làm gia
tăng nguy cơ quên lãng thảm họa (Mục~\ref{subsec:lang-adapt}). Vì vậy, chúng tôi chèn một
\textit{giai đoạn căn chỉnh đóng băng}: đóng băng toàn bộ phần thân Transformer và chỉ huấn
luyện ma trận nhúng đầu vào, ma trận đầu ra và độ chệch của nó trên một lượng nhỏ ngữ liệu
tiếng Việt với mục tiêu MLM.

Trực giác của bước này là để hai ma trận vừa khởi tạo \textit{căn chỉnh với nhau và với phần
thân cố định} trong không gian biểu diễn trước khi cho phép phần thân thay đổi; nhờ đó điểm
khởi đầu của giai đoạn tiếp tục tiền huấn luyện ổn định hơn. Cách làm này nằm trong tinh thần
chiến lược \textit{đóng băng rồi thích nghi} của de Vries và
Nissim~\cite{devries2021recycle}, vốn cho thấy việc tách giai đoạn học lớp nhúng khỏi giai
đoạn tinh chỉnh toàn bộ giúp tái sử dụng mô hình hiệu quả hơn. Trong thực nghiệm
(Chương~\ref{chap:experiment}), chúng tôi quan sát thấy giai đoạn này cải thiện chất lượng
khởi đầu, do đó được giữ lại trong quy trình.

% ======================================================================
\section{Tiếp tục tiền huấn luyện thích nghi}
\label{sec:cpt}

Ở giai đoạn cuối, toàn bộ mô hình được tinh chỉnh trên ngữ liệu tiếng Việt quy mô lớn
CulturaX với mục tiêu MLM. Theo Gururangan và cộng sự~\cite{gururangan2020dapt}, việc tiếp
tục tiền huấn luyện trên dữ liệu miền đích mang lại cải thiện ổn định; ở đây miền đích là
toàn bộ ngôn ngữ tiếng Việt.

\textbf{Lịch tốc độ học WSD.} Chúng tôi sử dụng lịch \textit{Warmup--Stable--Decay}
(WSD)~\cite{hu2024minicpm} thay cho lịch cosine cố định, vì WSD cho phép huấn luyện theo
nhiều phiên nối tiếp mà không phải cam kết trước tổng số bước. Tốc độ học $\eta$ theo bước
toàn cục $s$ được định nghĩa như Công thức~\eqref{eq:wsd}:
\begin{equation}
  \label{eq:wsd}
  \eta(s) =
  \begin{cases}
    \eta_{\mathrm{peak}} \cdot \dfrac{s}{S_w} & \text{nếu } s < S_w \quad \text{(khởi động)}\\[2mm]
    \eta_{\mathrm{peak}} & \text{nếu } S_w \le s < S_c \quad \text{(ổn định)}\\[2mm]
    \eta_{\mathrm{peak}} \cdot \Big(1 - \sqrt{\tfrac{\,s - S_c\,}{S_d}}\,\Big) & \text{nếu } s \ge S_c \quad \text{(làm nguội)}
  \end{cases}
\end{equation}
Trong đó $S_w$ là mốc kết thúc khởi động, $S_c$ là mốc bắt đầu làm nguội, và $S_d$ là độ dài
giai đoạn làm nguội. Pha \textit{ổn định} giữ tốc độ học ở đỉnh trong phần lớn quá trình; pha
\textit{làm nguội} dùng dạng $1-\sqrt{\cdot}$ vốn được Hägele và cộng sự~\cite{hagele2024scaling}
chứng minh là dạng làm nguội cho chất lượng hạ nguồn tốt nhất ở một ngân sách cho trước. Vì
lịch WSD được điều khiển theo \textit{bước toàn cục}, một phiên huấn luyện nối tiếp tiếp tục ở
pha ổn định mà không bị khởi động lại đột ngột.

\textbf{Khởi động lại tốc độ học.} Một điểm kiểm tra đã được làm nguội về tốc độ học gần 0
gần như không thể cập nhật trọng số nếu huấn luyện tiếp mà không xử lý. Do đó, mỗi khi tiếp
tục huấn luyện trên dữ liệu mới (ví dụ một pha làm nguội trên dữ liệu mới), tốc độ học được
\textit{khởi động lại} rồi làm nguội lại theo đúng khuyến nghị của Ibrahim và cộng
sự~\cite{ibrahim2024continual}. Ngoài ra, bộ tối ưu AdamW được giữ \textit{liên tục} giữa các
phiên (mang theo các mô-men bậc nhất và bậc hai) với cặp tham số $\beta = (0{,}9;\, 0{,}95)$
theo công thức gốc của NeoBERT, để bộ tiền điều kiện của bộ tối ưu khớp với chế độ mà phần
thân đã được huấn luyện.

% ======================================================================
\section{Tổng kết phương pháp}
\label{sec:method-summary}

Tóm lại, phương pháp xây dựng ViNeoBERT gồm ba giai đoạn: (1)~khởi tạo lớp nhúng và lớp đầu
ra bằng phép chiếu SALT theo từng token từ ViDeBERTa sang không gian NeoBERT, với tập neo
được làm giàu bằng cặp dịch và độ chệch đầu ra khởi tạo theo tần suất; (2)~căn chỉnh đóng
băng để hai ma trận mới ổn định cùng phần thân; và (3)~tiếp tục tiền huấn luyện thích nghi
theo lịch WSD. So với SALT gốc, đề tài có ba khác biệt chính: \textit{mở rộng tập neo bằng
cặp dịch} (Mục~\ref{sec:anchor}), \textit{áp dụng phép chiếu cho cả lớp đầu ra tách rời kèm
khởi tạo độ chệch theo tần suất} (Mục~\ref{sec:salt-proj}), và \textit{bổ sung giai đoạn căn
chỉnh đóng băng} (Mục~\ref{sec:freeze}). Hiệu quả của các thành phần này được đánh giá trong
Chương~\ref{chap:experiment}.
